% ===========================================================================
%  Writeup 20261119-142748
%
%  A write-up: one problem environment per question, and beneath it a marked
%  solution region holding the agreed argument. `board writeup add <label>`
%  writes both and rebuilds the PDF; a region stays empty until its answer is
%  agreed. Code goes in verbatim, so nothing here needs shell-escape.
% ===========================================================================
\documentclass[11pt]{article}
% The subject's own preamble when it has one (`latex/coursemacros.sty`, on
% TEXINPUTS through `board build`), else a plain one.
\IfFileExists{coursemacros.sty}{\usepackage{coursemacros}}{%
  \usepackage{amsmath}\usepackage{amssymb}\usepackage{amsthm}%
  \usepackage[margin=1in]{geometry}\usepackage{graphicx}}
% The board's macro vocabulary, every one a \providecommand, so a card's
% notation typesets here too and the subject's definitions still win.
\IfFileExists{board-macros.tex}{\input{board-macros.tex}}{}
\ifdefined\problem\else
\newenvironment{problem}[1]
  {\par\medskip\noindent\textbf{Problem #1.}\ \itshape}
  {\par\medskip}
\fi
\ifdefined\solution\else
\newenvironment{solution}
  {\par\noindent\textbf{Solution.}\ }
  {\par\medskip}
\fi
\providecommand{\todo}[1]{\textbf{[TODO: #1]}}

\title{Writeup 20261119-142748}
\author{}
\date{\today}

\begin{document}
\maketitle

% One problem/solution pair per question, in the order the sheet or the
% session posed them.

\begin{problem}{append-cost}
Under the doubling rule the grow fires at counts $1, 3, 7, 15, \dots$ and copies whatever is already there each time. Add up the elements copied across every grow during $n$ appends, and say how that total grows as $n$ grows.
\end{problem}
% ===== SOLUTION append-cost =====
Accounting: charge 3 units per element.

- 1 unit for insertion.
- 1 unit for when it gets copied at the next array double.
- 1 unit for copying its $k$-offset twin.

The array is two halves of $n/2$: a $k$-offset element in the new half, and its twin at the same offset in the old half.

So at the next doubling, everything is paid for.

$$O(3n) = O(n)$$
% ===== END SOLUTION append-cost =====

\end{document}
